\begin{satz}{Satz von Milman--Schechtman}{SatzVonMilmanSchechtman}
Es gibt ein $C_{BS} \in [1, \infty[$ derart, dass für alle $m, n \in \N$
mit $\frac 1 {10} \log(n) \le m \le n$ und jeden $n$--dimensionalen normierten
$\R$--Vektorraum $(E, \norm\cdot_E)$ ein $F \le E$ existiert mit $\dim F = m$
und
\[d_{BM}\left((F, \restrict{\norm\cdot_E}F), \ell_2^m \right)
\le 3 \frac{\sqrt{2 \log\left(\frac {15}2 \right)}}
{\sqrt{\log\left(\frac {15}2 \right)} - 1}C_{BS}
\sqrt{\frac m {\log\left(1+ \frac n m \right)}}\text.\]
\end{satz}
\begin{proof}
Nach \cite{BS} (dort: Theorem 2) gibt es ein $C_{BS} \in \R_{>0}$, sodass für
alle normierten Räume $(X, \norm\cdot_X)$ ein $M \in \N$ mit
$M > \frac 1 2 \dim X$ und $y \in X^M$ existieren mit
$\norm\cdot_{\infty,M}\le\norm\cdot_X\circ\bprod\cdot y_M
\le C_{BS}\norm\cdot_{2,M}$.\\
Wegen $\norm{e_1^M}_{\infty, M} = \norm{e_1^M}_{2, N}$ ist $C_{BS} \ge 1$.\\
Seien also $m, n \in \N$ mit $\frac 1 {10} \log(n) \le m \le n$ und
$(E, \norm\cdot_E)$ ein $n$--dimensionaler normierter $\R$--Vektorraum.\\
Nach dem oben zitierten Satz gibt es also ein $\widetilde M \in \N$ mit
$\widetilde M > \frac n 2$ und ein $\widetilde y \in E^{\widetilde M}$ mit
$\norm\cdot_{\infty,\widetilde M}
\le\norm\cdot_E\circ\bprod\cdot {\widetilde y}_{\widetilde M}
\le C_{BS}\norm\cdot_{2,\widetilde M}$.\\
Sei $N \coloneqq  \floor{\frac n 2}$. Dann ist $N \le \frac n 2 < \widetilde M$ und
daher $x\coloneqq \left(\frac 1{C_{BS}}\widetilde y_i \right)_{i=1}^N$ wohldefiniert
und es gilt (durch Auffüllung mit Nullen auf den Indizes ab $N+1$):
\begin{align}\frac 1{C_{BS}}
\norm\cdot_{\infty, N} \le \norm\cdot_E \circ \bprod\cdot x_N \le
\norm\cdot_{2,N} \text.\label{svmseq:1}\end{align}
Damit ist $x \neq 0$, sogar $x_i \neq 0$, da
$\frac 1 {C_{BS}} = \frac 1 {C_{BS}} \norm{e_i^N}_{\infty,N} \le \norm{x_i}_E$
für alle $i \in \haken N$.\\
Seien nun $T\coloneqq S^{m-1}$,
\[\bigg((\Omega, \mfa, \Prim), (\gamma,\delta,\widetilde g),(N, m, mN)\bigg)
\text{ ein }(N(0,1), 3)\text{--Tupel zu }(3,(N,m,mN))\] (existiert nach
\ref{ExistenzNormalNTupel}), $G_x \coloneqq  \bprod \cdot x_N \circ g$ und $(Z_{\gamma, x}, M_{\delta,T})$ das
Auswertungspaar zu $(\gamma, \delta, x, T)$.\\
Seien weiter
\[f: = \inf_{t \in T} \, \norm\cdot_E \circ \bprod \cdot x _N \circ G_x,
\quad
h\coloneqq  \sup_{t \in T} \, \norm\cdot_E \circ \bprod \cdot x _N \circ G_x\text.\]
Nach \ref{IntegrierbarkeitVonExtremaInGordonschenPaaren} sind
\begin{align}
f, h\text{ integrierbar und es ist } f \le h\text.\label{svmseq:2}
\end{align}
\underline{1. Fall:} $m \le \frac 2 {13}(N+1)$\absatz
\underline{1.1. Fall:} $\E(Z_{\gamma, x}) < 2a_m$\absatz
Seien $\mu_{j, N}$ für alle $j \in \haken N$ gemäß \ref{ktesMaximum},
\[\norm\cdot_{(m)}:\R^N \rightarrow [0, \infty[, v \mapsto
\frac 1 m \sum_{i=1}^m \mu_{i, N}\left(\pfeil {\abs\cdot}N (v)\right)\]
und
\[ V\coloneqq  \mengea{z \in \R^N\setminus\meng 0}
    {\left(\forall\, i \in \haken N: \abs{z_i} \in\meng{0,\nicefrac 1m} \right)
    \wedge \left(\abs{\menge{i \in \haken N}{\abs{z_i} \neq 0}}\le m \right)}
    \text.\]
Dann ist $V \subseteq \R^N$, nichtleer,
beschränkt (da endlich) und es ist $0 \notin V$. \hfill $(a)$\\
Nach \refclaim{Sortierungsnorm}{Sortierungsnorm3} ist
\begin{align}
\norm\cdot_{(m)}\text{ eine
Norm auf dem }\R^N\text{ mit } 
\norm\cdot_{(m)} \le \norm\cdot_{\infty, N}\text.\label{svmseq:3}
\end{align}
Ferner ist nach \refclaim{Sortierungsnorm}{Sortierungsnorm1}
$\sup\menge{\norm v_{2, N}}{v \in V} = \frac 1{\sqrt m}$ \hfill $(b)$\\
und nach \refclaim{Sortierungsnorm}{Sortierungsnorm2}
\begin{align}
\forall\, w \in \R^N: \sup_{v \in V}\sprod w v_N = \norm w_{(m)}\text.\tag{$c$}
\end{align}
Für alle $s \in \R^m$ und alle $\omega \in \Omega$ gilt:
\begin{align}
&\,(\bprod \cdot x \circ ev(s, \cdot) \circ g)(\omega)
= \bprod {g(\omega)s}x_N = (\bprod \cdot x_N \circ g(\omega))(s)\notag\\
=&\, ev(s, \cdot) \bigg(\bprod \cdot x _N \circ g(\omega) \bigg)
= (ev(s, \cdot) \circ \bprod \cdot x _N \circ g)(\omega) \text.\label{svmseq:4}
\end{align}
Wegen $N = \floor {\frac n 2}$ ist $N+1\ge \frac n 2$, also auch
$\frac {N+1}m \ge \frac n {2m}$ und damit
\begin{align}\left(1+\frac{N+1}m \right)^2
\ge 1 + 2 \frac{N+1}m \ge 1 + \frac n m\, , \text{ also } 1 + \frac{N+1}m
\ge \sqrt{1+ \frac n m} \text.\label{svmseq:5}\end{align}
Wegen $m \le \frac 2 {13}(N+1)$ und damit
$1+\frac{N+1}m \ge 1+\frac{13}2 = \frac{15}2$ gilt weiter:
\begin{align}
&\,\E\left(\norm \cdot_{(m)} \circ \gamma \right) - \frac {a_m}{\sqrt m}
\ge \E\left(\gamma_m^\ast \right) - \frac {a_m}{\sqrt m}
\overset{\refclaim{AbschaetzungEinesGammafunktionsquotienten}
{AbschaetzungEinesGammafunktionsquotienten2}}\ge 
\E\left(\gamma_m^\ast \right)-1\notag\\
\overset{\refclaim
{UntereSchrankeDesErwartungswertsVonAbsolutenN01Ordnungsstatistiken}
{UntereSchrankeDesErwartungswertsVonAbsolutenN01Ordnungsstatistiken3}}\ge
&\,
\sqrt{\log\left(1 + \frac {N+1}m \right)}-1\notag\\
=&\, \sqrt{\log\left(1 + \frac {N+1}m \right)} - 
\left(1 - \frac{\sqrt{\log\left(\frac {15}2 \right)} -1}
{\sqrt{\log\left(\frac {15}2 \right)}} \right)
\sqrt{\log\left(\frac {15}2 \right)}\notag\\
\ge &\,\sqrt{\log\left(1 + \frac {N+1}m \right)} - 
\left(1 - \frac{\sqrt{\log\left(\frac {15}2 \right)} -1}
{\sqrt{\log\left(\frac {15}2 \right)}} \right)
\sqrt{\log\left(1 + \frac {N+1}m \right)}\notag\\
=&\, \left(\frac{\sqrt{\log\left(\frac {15}2 \right)} -1}
{\sqrt{\log\left(\frac {15}2 \right)}} \right)
\sqrt{\log\left(1 + \frac {N+1}m \right)}\notag\\
\overset{\eqref{svmseq:5}}\ge 
&\, \left(\frac{\sqrt{\log\left(\frac {15}2 \right)} -1}
{\sqrt{\log\left(\frac {15}2 \right)}} \right)
\sqrt{\frac 1 2\log\left(1 + \frac {n}m \right)}\notag\\
=& \underbrace{\left(\frac{\sqrt{\log\left(\frac {15}2 \right)} -1}
{\sqrt{2\log\left(\frac {15}2 \right)}} \right)}_{>0, \text{ da }
\frac{15}2 > \exp(1)}
\sqrt{\log\left(1 + \frac {n}m \right)} \eqqcolon \alpha \text.\label{svmseq:6}
\end{align}
Seien $E_1\coloneqq  \R^N$, $\norm\cdot_{E_1}\coloneqq  \norm\cdot_{(m)}$,
$x_1 \coloneqq  (e_i^N)_{i=1}^N$ und $G_{x_1}\coloneqq  \bprod \cdot {x_1}_N \circ g = g$.\\
Wegen $(a), (b), (c)$ sind die Voraussetzungen von
\ref{GordonschePaareUndNormiterierendeMengen} für $(E_1, \norm\cdot_{E_1})$,
$x_1$, $T$, $V$ und $c\coloneqq \frac 1 {\sqrt{m}}$ erfüllt und damit gilt:
\begin{align}
\E\left(\norm\cdot_{(m)} \circ \gamma \right) - \frac {a_m}{\sqrt m}
& \overset{\ref{ErwartungswertDer2NormEinesStandardnormalverteiltenProzesses}}
= \E\left(\norm\cdot_{(m)} \circ \bprod \cdot {x_1}_N \circ\gamma \right)
- \frac 1 {\sqrt m} \E\left(\norm\cdot_{2, m} \circ \delta \right)\notag\\
&\overset{\ref{IterierteEuklidischeNorm}}=
\E\left(\norm\cdot_{(m)} \circ \bprod \cdot {x_1}_N \circ\gamma \right)
+ \frac 1 {\sqrt m}\E\left(\inf_{t \in T} \bprod t \delta_m \right)\notag\\
&\overset{\refclaim{GordonschePaareUndNormiterierendeMengen}
{GordonschePaareUndNormiterierendeMengen1}}\le
\E\left(\inf_{t \in T}\, \norm\cdot_{(m)}\circ ev(t,\cdot)\circ G_{x_1} \right)
\text.\label{svmseq:7}
\end{align}
Außerdem sind wegen $x \neq 0$ und \eqref{svmseq:1} auch die Voraussetzungen
von \ref{GordonschePaareUndSubeuklidischeBildhalbnormen} für
$(E, \norm\cdot_E)$, $x$ und $T$ erfüllt, sodass gilt:
\begin{align}
\E\left(\sup_{t \in T}\, \norm\cdot_{E}\circ ev(t,\cdot)\circ G_{x} \right)
&\overset{
\refclaim{GordonschePaareUndSubeuklidischeBildhalbnormen}
{GordonschePaareUndSubeuklidischeBildhalbnormen1}}
\le \E\left(\norm\cdot_E \circ \bprod \cdot x_N \circ \gamma \right)
+ \E\left(\sup_{t \in T} \bprod t \delta_m \right)\notag\\
&\overset{\ref{IterierteEuklidischeNorm}}=
\E\left(\norm\cdot_E \circ \bprod \cdot x_N \circ \gamma \right)
+ \E\left(\norm\cdot_{2, m}\circ \delta \right)\notag\\
&\overset{\ref{ErwartungswertDer2NormEinesStandardnormalverteiltenProzesses}}=
\E\left(\norm\cdot_E \circ \bprod \cdot x_N \circ \gamma \right) + a_m\notag\\
&\overset{\text{1.1 Fall}}< 3a_m \text.\label{svmseq:8}
\end{align}
Wir erhalten also insgesamt:
\begin{align*}
0 & \overset{\eqref{svmseq:6}}<
\frac 1 {C_{BS}}\left(\frac{\sqrt{\log\left(\frac {15}2 \right)} -1}
{\sqrt{2\log\left(\frac {15}2 \right)}} \right)
\sqrt{\log\left(1 + \frac {n}m \right)}\\
&\overset{\eqref{svmseq:6}}\le 
\frac 1 {C_{BS}}\left(\E\left(\norm \cdot_{(m)} \circ \gamma \right) 
- \frac {a_m}{\sqrt m}\right)\\
&\overset{\eqref{svmseq:7}}\le
\E\left(\inf_{t \in T}\, 
\frac 1 {C_{BS}}\norm\cdot_{(m)}\circ ev(t,\cdot)\circ G_{x_1} \right)\\
&=
\E\left(\inf_{t \in T}\, 
\frac 1 {C_{BS}}\norm\cdot_{(m)}\circ ev(t,\cdot)\circ g \right)\\
&\overset{\eqref{svmseq:3}}\le
\E\left(\inf_{t \in T}\, 
\frac 1 {C_{BS}}\norm\cdot_{\infty, N}\circ ev(t,\cdot)\circ g \right)\\
&\overset{\eqref{svmseq:1}}\le
\E\left(\inf_{t \in T}\, 
\norm\cdot_{E}\circ \bprod \cdot x_N \circ ev(t,\cdot)\circ g \right)\\
&\overset{\eqref{svmseq:4}}=
\E\left(\inf_{t \in T}\, 
\norm\cdot_{E}\circ ev(t,\cdot)\circ \bprod \cdot x_N \circ g \right)
=\E\left(\inf_{t \in T}\, \norm\cdot_{E}\circ ev(t,\cdot)\circ G_x \right)\\
&= \E(f) \overset{\eqref{svmseq:2}}\le \E(h)
= \E \left(\sup_{t \in T}\, \norm\cdot_{E}\circ ev(t,\cdot)\circ G_x \right)
\overset{\eqref{svmseq:8}}< 3 a_m\text.
\end{align*}
Da $f \ge 0, h \ge 0$ und beide integrierbar sind, gibt es nach
\ref{Erwartungswertlemma} ein $\omega_0 \in \Omega$ mit
$f(\omega_0)>0$ und $\frac{h(\omega_0)}{f(\omega_0)} \le \frac{3a_m}\alpha$.\\
Sei $A\coloneqq  G_x(\omega_0)$. Dann ist $A \in L(\ell_2^m, E)$ und es gilt:
\[\inf_{t \in S^{m-1}}\norm{At}_E =
  \inf_{t \in T} \norm{G_x(\omega_0)t}_E =
  \left(\inf_{t \in T} \norm\cdot_E \circ ev(t, \cdot) \circ G_x\right)
  (\omega_0) = f(\omega_0) > 0\]
und damit ist $A$ nach
\ref{InjektivitätLinearerAbbildungenZwischenNormiertenRäumen} injektiv.\\
Sei $F\coloneqq  \Bild A$. Da $A$ injektiv ist, ist $\dim F = m$ und da jeder
endlichdimensionale normierte Raum vollständig ist, gilt:
\begin{align*}
&\,d_{BM}\bigg((F, \norm\cdot_E\vert_F), \ell_2^m \bigg)
\le \opnorm A {\ell_2^m}{F} \opnorm{A\inv}{F}{\ell_2^m}
\overset{\ref{OperatornormDerInversen}}=
\frac{\Sup{t \in T} \norm{At}_E}{\Inf{t \in T}\norm{At}_E}\\
=&\, \frac{\Sup{t \in T}\norm{G_x(\omega_0)t}_E}
       {\Inf{t \in T}\norm{G_x(\omega_0)t}_E}
= \frac{\left(\Sup{t \in T} \norm\cdot_E \circ ev(t, \cdot) \circ G_x\right)
  (\omega_0)}
  {\left(\Inf{t \in T} \norm\cdot_E \circ ev(t, \cdot) \circ G_x\right)
  (\omega_0)} = \frac{h(\omega_0)}{f(\omega_0)} \le \frac{3a_m}\alpha\\
\overset{C_{BS} \ge 1}\le&\, \frac{3a_m}
{\frac 1 {C_{BS}}\left(\frac{\sqrt{\log\left(\frac {15}2 \right)} -1}
{\sqrt{2\log\left(\frac {15}2 \right)}} \right)
\sqrt{\log\left(1 + \frac {n}m \right)}}
= 3 \frac{\sqrt{2 \log\left(\frac {15}2 \right)}}
{\sqrt{\log\left(\frac {15}2 \right)} - 1}C_{BS}
\frac{a_m}{\sqrt{\log\left(1 + \frac {n}m \right)}}\\
\overset{\refclaim{AbschaetzungEinesGammafunktionsquotienten}
{AbschaetzungEinesGammafunktionsquotienten2}}\le &\,
3 \frac{\sqrt{2 \log\left(\frac {15}2 \right)}}
{\sqrt{\log\left(\frac {15}2 \right)} - 1}C_{BS}
\sqrt{\frac m {\log\left(1 + \frac {n}m \right)}}\text.
\end{align*}
Damit ist der Fall dieser Fall abgeschlossen.\absatz
\underline{1.2. Fall:} $\E(Z_{\gamma, x}) \ge 2 a_m$\absatz
Sei dann zusätzlich $\eps \coloneqq  \frac 1 2$.\\
Wegen \eqref{svmseq:1} ist $x \neq 0$ und 
$\norm\cdot_E\circ \bprod \cdot x_N \le \norm\cdot_{2,N}$. \hfill $(a')$\\
Ferner ist $T\neq \emptyset$, $T \subseteq S^{m-1}$ und $T=-T$. \hfill $(b')$\\
Schließlich gilt:
\begin{align}\E(M_{\delta, T}) \overset{\ref{IterierteEuklidischeNorm}}=
 \E\left(\norm\cdot_{2,m} \circ \delta \right)
 \overset{\ref{ErwartungswertDer2NormEinesStandardnormalverteiltenProzesses}}=
 a_m = \frac 1 2 2a_m \le \eps\E(Z_{\gamma,x})\text. \tag{$c'$}
 \end{align}
Wegen $(a'), (b'), (c')$ sind die Voraussetzungen von
\ref{ExistenzEingeschraenkterFastIsometrien} erfüllt und es gibt ein 
$A \in L(\ell_2^m, E)$ injektiv mit
\begin{align}
\opnorm A {\ell_2^m} E \opnorm{A\inv}{A(E)}{\ell_2^m} \le
\frac{1+\eps}{1-\eps} = \frac{\nicefrac 3 2}{\nicefrac 1 2} = 3\text.
\label{svmseq:9}
\end{align}
Sei nun $\vi:\!\ ]1,\infty[\rightarrow\R,s\mapsto\frac{\log(s)}{\log(1+s)}$.\\
Dann ist $\vi$ offenbar stetig differenzierbar und es gilt für alle
$z\in\!\ ]1,\infty[$:
\[\vi'(z)=\frac{\frac{\log(1+z)}{z}-\frac{\log(z)}{1+z}}{\log(1+z)^2}>0\text,\]
also $\vi$ monoton wachsend, sodass wegen $n \ge 2$ und
$m \ge \frac 1{10}\log(n)$ gilt:
\begin{align}\frac m{\log\left(1+\frac n m \right)}
\ge \frac 1 {10}\frac{\log(n)}{\log(n+1)}
= \frac 1 {10}\vi(n) \ge \frac 1 {10}\vi(2) 
=\frac 1 {10}\frac{\log(2)}{\log(3)}\text. \label{svmseq:10}\end{align}
Ferner gilt:
\begin{align}\frac{\sqrt{2 \log\left(\frac {15}2 \right)}}
{\sqrt{\log\left(\frac {15}2 \right)} - 1}
\sqrt{\frac 1 {10}\frac{\log(2)}{\log(3)}} > 1\text. \label{svmseq:11}
\end{align}
Sei nun $F\coloneqq  \Bild A$. Dann gilt:
\begin{align*}
d_{BM}\bigg((F, \norm\cdot_E\vert_F), \ell_2^m \bigg)
&\le \opnorm A {\ell_2^m} F \opnorm{A\inv}{F}{\ell_2^m} 
\overset{\eqref{svmseq:9}}\le 3
\overset{C_{BS}\ge 1}\le 3C_{BS}\\
&\overset{\eqref{svmseq:11}}< 3\frac{\sqrt{2 \log\left(\frac {15}2 \right)}}
{\sqrt{\log\left(\frac {15}2 \right)} - 1}
C_{BS}\sqrt{\frac 1 {10}\frac{\log(2)}{\log(3)}}\\
&\overset{\eqref{svmseq:10}}\le
3\frac{\sqrt{2 \log\left(\frac {15}2 \right)}}
{\sqrt{\log\left(\frac {15}2 \right)} - 1}
C_{BS}\sqrt{\frac m{\log\left(1+\frac n m \right)}}\text.
\end{align*}
Damit ist auch dieser Fall abgeschlossen.\absatz
\underline{2. Fall:} $m > \frac 2 {13}(N+1)$\absatz
Wegen $N = \floor{\frac n 2} \ge \frac n 2 - 1$ ist 
$m > \frac 2 {13}(N+1)\ge \frac 2 {13}\cdot\frac n 2 = \frac n {13}$ und
damit $\log\left(1+\frac n m \right)< \log(14)$.\\
Ferner gilt
\[3\frac{\sqrt{2 \log\left(\frac {15}2 \right)}}
{\sqrt{\log\left(\frac {15}2 \right)} - 1} \frac 1 {\sqrt{\log(14)}}> 1\text.\]
Sei nun $F \le E$ mit $\dim F = m$.\\
Dann ist nach dem Satz von John
\begin{align*}
d_{BM}\left((F, \restrict{\norm\cdot_E}F), \ell_2^m \right)
&\le \sqrt m \overset{C_{BS}\ge 1}\le C_{BS}\sqrt{m}
< 3\frac{\sqrt{2 \log\left(\frac {15}2 \right)}}
{\sqrt{\log\left(\frac{15}2 \right)}-1}C_{BS}\frac{\sqrt m}{\sqrt{\log(14)}}\\
& \le  3\frac{\sqrt{2 \log\left(\frac {15}2 \right)}}
{\sqrt{\log\left(\frac {15}2 \right)} - 1}
C_{BS}\sqrt{\frac m{\log\left(1+\frac n m \right)}}\text.
\end{align*}
\end{proof}